% ---------------------------------------------------------------
% seccoes/terminologia.tex — Terminologia específica
% ---------------------------------------------------------------

\chapter*{Terminologia específica}
\addcontentsline{toc}{chapter}{Terminologia específica}

Termos técnicos usados neste relatório, por ordem alfabética.

\begin{tabela}{L{\dimexpr0.2182\linewidth-2\tabcolsep\relax}L{\dimexpr0.7798\linewidth-2\tabcolsep\relax}}
\caption{Terminologia específica}\label{tab:1}\\
\linhacab \cab{\textbf{Termo}} & \cab{\textbf{Definição}} \\
\endfirsthead
\linhacab \cab{\textbf{Termo}} & \cab{\textbf{Definição}} \\
\endhead
Arbitragem & Escolha, no fim de cada ciclo, da decisão de velocidade e da decisão de via entre todas as propostas das regras. \\
\rowcolor{corFundo}Base de conhecimento & Conjunto das regras que representam o conhecimento do domínio. \\
Canal & Cada uma das duas partes independentes da decisão: velocidade e via. \\
\rowcolor{corFundo}Encadeamento progressivo & Inferência que parte dos factos e aplica as regras cujas condições se verificam (\emph{forward chaining}). \\
Encadeamento regressivo & Inferência que parte de um objetivo e procura os factos que o provam (\emph{backward chaining}). \\
\rowcolor{corFundo}\emph{Entry point} & Em \emph{Drools}, fluxo com nome pelo qual os eventos de um tipo entram na sessão. \\
Evento & Leitura de um sensor num dado instante, que expira ao fim de algum tempo; distingue-se do facto, que descreve o estado atual. \\
\rowcolor{corFundo}Explicação \emph{Why / How / Why Not} & Explicações de porque foi tomada uma decisão, de como se chegou à mesma e de porque não ter sido tomada outra. \\
Facto & Unidade de informação sobre a situação atual, mantida na memória de trabalho. \\
\rowcolor{corFundo}Lógica difusa & Raciocínio com categorias, a que um valor pertence com um grau entre 0 e 1 (\emph{fuzzy logic}). \\
Memória de trabalho & Conjunto dos factos que o sistema conhece num dado momento (\emph{working memory}). \\
\rowcolor{corFundo}Metaconhecimento & Conhecimento sobre as próprias regras, que controla a inferência: ordem das fases, prioridades, desempate. \\
Motor de inferência & Componente que avalia as regras e produz conclusões. \\
\rowcolor{corFundo}Nível de risco & Gravidade de uma situação: baixo, moderado, alto ou crítico. Acompanha a decisão, mas não a seleciona. \\
Prioridade & Categoria, de 10 a 70, que decide entre propostas concorrentes. \\
\rowcolor{corFundo}Processamento de eventos complexos (CEP) & Deteção de padrões em sequências de eventos ao longo do tempo, como três leituras de TTC baixo em 2 s (\emph{complex event processing}). \\
Sistema pericial & Sistema que representa conhecimento especializado de forma explícita e o utiliza, através de um motor de inferência, para chegar a conclusões justificáveis (\emph{expert system}). \\
\rowcolor{corFundo}TTC & Tempo até à colisão (\emph{time to collision}): distância ao veículo da frente a dividir pela velocidade de aproximação. \\
\emph{WebSocket} & Protocolo de comunicação persistente e bidirecional sobre uma única ligação TCP (RFC 6455). \\
\end{tabela}
